\section{本章小结}
\label{sec:weakd-summary}

本章开头问：对只在几乎处处意义下确定、可能有尖点或跳跃的函数，能否用光滑函数逼近，能否给它一个可以运算的导数。第一个问题由局部平均回答。局部可积只要求函数在每个紧含于区域的集合上可积，这恰好使它与测试函数的乘积可以积分。以归一化的隆起函数为权重做局部平均，得到磨光函数；它的光滑性完全来自核，因为对参数求导可以全部移到核上，而局部可积的函数提供与参数无关的可积上界。

磨光函数逼近原函数的依据是平移的积分连续性。从外测度的矩形覆盖得到开集外逼近与紧集内逼近，再用矩形的有限并和简单函数逼近可积函数，就能证明小平移只引起小的积分误差。核非负且积分为$1$，使磨光误差成为平移误差的加权平均，所以在每个内部紧集上$\int_K\lvert f_\varepsilon-f\rvert\to0$。这种逼近只控制积分误差，阶跃函数的例子说明逐点最大误差可以不趋于零。由于磨光值本身就是一个测试积分，局部逼近立即给出测试函数判别：与全部测试函数积分为零的局部可积函数几乎处处为零。

第二个问题由测试恒等式回答。光滑函数的分部积分把导数转到测试函数上，测试函数的紧支撑使边界项消失，得到的右端只要求函数局部可积。弱导数就是能代替经典导数放进这个恒等式的局部可积函数；测试函数判别保证它几乎处处唯一，对光滑函数它与经典导数一致。

尖点与跳跃的对照说明了定义的边界：$\lvert x\rvert$在原点两侧的端点项相消，弱导数是$\operatorname{sgn}x$；阶跃函数留下读取原点值的规则，它是分布导数，却不能由任何局部可积函数表示。分布导数对每个局部可积函数都存在，弱导数还要求这条线性规则由局部可积函数给出。

在一维有界区间上，函数及其弱导数可积，等价于函数有绝对连续代表并能由弱导数积分恢复；弱导数为零的函数是常数，这一点依靠积分为零的测试函数恰好是测试函数的导数。端点值属于唯一的连续代表，局部可积的情形只在内部紧区间上得到绝对连续代表。最后，磨光与弱导数可以交换次序，因此在每个内部紧集上，光滑的磨光函数及其经典导数同时在积分意义下逼近原函数及其弱导数。

这些结论都是局部的，并且用积分误差衡量。把函数与若干阶弱导数的$L^p$大小合成范数，研究空间的完备性、整个区域上的光滑稠密性、边界处理与嵌入，是第~\ref{chap:sobolev-spaces}章的内容；磨光在$L^p$范数下的收敛由第~\ref{chap:function-approximation}章处理。下一章回到测度本身，研究一个测度何时能写成另一个测度的加权积分。
